\input{estilo}
% Propuesta: la misma tabla antes y después de un receive, a las 10:00:00.
% Visible es "visibility_timeout en el pasado". El receive toma la primera
% visible por id (m₁) y le corre el timeout 60 segundos al futuro; m₂ ya
% estaba tomado y sigue invisible. Las columnas llevan los nombres del SQL del
% texto. Agregado: las filas, las horas y la columna de estado.
\begin{document}
\begin{figura}\begin{tikzpicture}
  \def\h{0.62}
  % una tabla: \tabla{x}{título}{hora de m₁}{estado de m₁}{color de m₁}{resaltar}
  \newcommand\tabla[6]{\begin{scope}[xshift=#1cm]
    \node[anchor=south west] at (0,0.1) {#2};
    \foreach \x/\t in {0.4/id, 1.6/mensaje, 3.7/visibility\_timeout} \node[cod, text=suave] at (\x,-0.3) {\t};
    \draw (0,-0.6) -- (5.9,-0.6);
    \foreach \f/\id/\m/\t/\e/\c in {1/1/1/#3/#4/#5, 2/2/2/10:00:40/invisible/suave, 3/3/3/09:59:58/visible/segundo} {
      \node at (0.4,{-0.6-\h*(\f-0.5)}) {\id};
      \node at (1.6,{-0.6-\h*(\f-0.5)}) {\var{m}\textsubscript{\m}};
      \node[cod] at (3.7,{-0.6-\h*(\f-0.5)}) {\t};
      \node[etiq, text=\c] at (5.3,{-0.6-\h*(\f-0.5)}) {\e};
    }
    #6
    \draw (0,{-0.6-3*\h}) -- (5.9,{-0.6-3*\h});
  \end{scope}}

  \tabla{0}{antes}{09:59:30}{visible}{segundo}{}
  \tabla{7.6}{después del receive}{10:01:00}{invisible}{suave}
    {\draw[acento, rounded corners=2pt] (-0.1,-0.65) rectangle (6.0,{-0.6-\h+0.05});}
  \draw[implica] (6.15,-1.53) -- (7.35,-1.53);
  \node[etiq, text=suave, anchor=north] at (6.75,-1.75) {10:00:00};
  \node[etiq, text=suave, anchor=north west] at (0,-2.75) {visible = \texttt{visibility\_timeout < now()}};
  \node[etiq, text=acento, anchor=north west] at (7.6,-2.75) {el receive lo corre a \texttt{now() + 60 s}};
\end{tikzpicture}\end{figura}
\end{document}
